% GENERATED FILE. Edit canonical lessons, then run npm run generate:books.
% canonical-source-hash: fnv1a64:261a3640172b0a5b
% canonical-lessons: MR-C287-avlambun-ahe, MR-C287-gelo-asto, MR-W287-six-lines, MR-R287-depends-and-would

\chapter{It Depends, and What Would Have Been}
\label{ch:mr-it-depends-and-what-would-have-been}
\hlchaptermodality{287}

\begin{chapteropening}
\textbf{By the end of this chapter:} \emph{I can say what something depends on (… -var avlambūn āhe), say what would have happened (gelo asto), and write six connected lines about a trip.}

Complete MR-R287-depends-and-would: say what something depends on and what would have been, after writing six connected lines about a trip.
\end{chapteropening}

\section[\texorpdfstring{\mr{हवामानावर} \mr{अवलंबून} \mr{आहे} (havāmānāvar …)}{हवामानावर अवलंबून आहे (havāmānāvar …)}]{\mr{हवामानावर} \mr{अवलंबून} \mr{आहे} (havāmānāvar avlambūn āhe) — it depends}

\label{lesson:MR-C287-avlambun-ahe}



\begin{quote}
Say \textquotedblleft{}to depend\textquotedblright{}. (\emph{avlambūn asṇe}.) Then \textquotedblleft{}if it rains, I will stay at home\textquotedblright{}. (\emph{pāūs paḍlā tar mī gharī rāhīn}.)
\end{quote}

\subsection*{You'll want to know: … -\mr{वर} \mr{अवलंबून} \mr{आहे}}
\begin{quote}
\textbf{\mr{आम्ही} \mr{उद्या} \mr{जाऊ} \mr{का}? \mr{ते} \mr{हवामानावर} \mr{अवलंबून} \mr{आहे}.}
\end{quote}

\begin{quote}
\emph{āmhī udyā jāū kā? te havāmānāvar avlambūn āhe.} — \textquotedblleft{}Will we go tomorrow? It depends on the weather.\textquotedblright{}
\end{quote}

What it depends on takes \textbf{-\mr{वर}}, \textquotedblleft{}on\textquotedblright{}: \textbf{\mr{हवामानावर}} (\emph{havāmānāvar}), \textbf{\mr{तुमच्यावर}} (\emph{tumcyāvar}), \textquotedblleft{}on you\textquotedblright{}. The verb is \textbf{\mr{असणे}}, so the tense goes there: \textbf{\mr{ते} \mr{हवामानावर} \mr{अवलंबून} \mr{असेल}} (\emph{te havāmānāvar avlambūn asel}), \textquotedblleft{}it will depend on the weather\textquotedblright{}. To ask: with \textbf{\mr{कशावर}} (\emph{kaśāvar}, \textquotedblleft{}on what\textquotedblright{}): \textbf{\mr{कशावर} \mr{अवलंबून} \mr{आहे}?} (\emph{kaśāvar avlambūn āhe?})

\subsection*{Your turn}
Say these aloud:

\begin{itemize}
\raggedright
  \item \textquotedblleft{}It depends on the weather\textquotedblright{} — \emph{te havāmānāvar avlambūn āhe}
  \item \textquotedblleft{}It depends on you\textquotedblright{} — \emph{te tumcyāvar avlambūn āhe}
  \item \textquotedblleft{}What does it depend on?\textquotedblright{} — \emph{kaśāvar avlambūn āhe?}
\end{itemize}

\begin{tcolorbox}[breakable,colback=teal!4,colframe=teal!35!black,title={Before you move on}]
\textquotedblleft{}It depends on the weather\textquotedblright{}? (\textbf{\mr{ते} \mr{हवामानावर} \mr{अवलंबून} \mr{आहे}}.) \textquotedblleft{}On you\textquotedblright{}? (\textbf{\mr{तुमच्यावर}}.) \textquotedblleft{}What does it depend on?\textquotedblright{} (\textbf{\mr{कशावर} \mr{अवलंबून} \mr{आहे}?})
\end{tcolorbox}
\section[\texorpdfstring{\mr{मी} \mr{गेलो} \mr{असतो} (mī gelo asto)}{मी गेलो असतो (mī gelo asto)}]{\mr{मी} \mr{गेलो} \mr{असतो} (mī gelo asto) — what would have been}

\label{lesson:MR-C287-gelo-asto}



\begin{quote}
Say \textquotedblleft{}it depends on the weather\textquotedblright{}. (\emph{te havāmānāvar avlambūn āhe}.) Then \textquotedblleft{}If it rains tomorrow, I will stay at home\textquotedblright{}. (\emph{jar udyā pāūs paḍlā, tar mī gharī rāhīn}.)

That looks ahead at what may still happen. This lesson looks back at what did not.
\end{quote}

\subsection*{You'll want to know: -\mr{लो} \mr{असतो}, 'would have'}
\begin{quote}
\textbf{\mr{मी} \mr{गेलो} \mr{असतो}.}
\end{quote}

\begin{quote}
\emph{mī gelo asto.} — \textquotedblleft{}I would have gone.\textquotedblright{}
\end{quote}

The past, \textbf{\mr{गेलो}}, then \textbf{\mr{असतो}} (\emph{asto}), which takes the same endings: \textbf{\mr{तो} \mr{गेला} \mr{असता}} (\emph{to gelā astā}), \textbf{\mr{ती} \mr{गेली} \mr{असती}} (\emph{tī gelī astī}), \textbf{\mr{आम्ही} \mr{आलो} \mr{असतो}} (\emph{āmhī ālo asto}). A woman says \textbf{\mr{मी} \mr{गेले} \mr{असते}} (\emph{mī gele aste}).

\subsection*{You'll want to know: the if that did not happen}
\begin{quote}
\textbf{\mr{जर} \mr{मला} \mr{वेळ} \mr{असता}, \mr{तर} \mr{मी} \mr{आलो} \mr{असतो}.}
\end{quote}

\begin{quote}
\emph{jar malā veḷ astā, tar mī ālo asto.} — \textquotedblleft{}If I had had time, I would have come.\textquotedblright{}
\end{quote}

In the condition, \textbf{\mr{असता}} (\emph{astā}) alone says \textquotedblleft{}had been\textquotedblright{}; with a verb, the past comes first: \textbf{\mr{जर} \mr{पाऊस} \mr{पडला} \mr{नसता}} (\emph{jar pāūs paḍlā nastā}), \textquotedblleft{}if it had not rained\textquotedblright{}, with \textbf{\mr{नसता}} (\emph{nastā}, \textquotedblleft{}had not been\textquotedblright{}). Both halves wear \textbf{\mr{असता}}, and that is how a listener knows it did not happen: \textbf{\mr{जर} \mr{तू} \mr{आला} \mr{असतास}, \mr{तर} \mr{आम्ही} \mr{बोललो} \mr{असतो}} (\emph{jar tū ālā astās, tar āmhī bollo asto}), \textquotedblleft{}if you had come, we would have talked\textquotedblright{}.

\subsection*{Your turn}
Say these aloud:

\begin{itemize}
\raggedright
  \item \textquotedblleft{}I would have gone\textquotedblright{} — \emph{mī gelo asto}
  \item \textquotedblleft{}If I had had time, I would have come\textquotedblright{} — \emph{jar malā veḷ astā, tar mī ālo asto}
  \item \textquotedblleft{}If it had not rained, we would have gone\textquotedblright{} — \emph{jar pāūs paḍlā nastā, tar āmhī gelo asto}
\end{itemize}

\begin{tcolorbox}[breakable,colback=teal!4,colframe=teal!35!black,title={Before you move on}]
\textquotedblleft{}I would have gone\textquotedblright{}? (\textbf{\mr{मी} \mr{गेलो} \mr{असतो}}.) \textquotedblleft{}She would have come\textquotedblright{}? (\textbf{\mr{ती} \mr{आली} \mr{असती}}.) \textquotedblleft{}If I had had time\textquotedblright{}? (\textbf{\mr{जर} \mr{मला} \mr{वेळ} \mr{असता}}.) How do you know it did not happen? (Both halves wear \textbf{\mr{असता}}.)
\end{tcolorbox}
\section[(six lines)]{Six lines about a trip}

\label{lesson:MR-W287-six-lines}



\begin{quote}
The A1 message asked for a few practised lines to Meera. This task asks for six lines, and they tell a story: a trip, to a village or anywhere else.

Pick one reader: Meera, or a friend, by name. Say \textbf{\mr{ऋतू}} and \textbf{\mr{ज्ञान}} aloud; both will go in.
\end{quote}

\subsection*{What to know first}
A told story is not a list of events. It has a \textbf{scene}, in \textbf{-\mr{त} \mr{होतो}} or \textbf{\mr{असताना}}. It has \textbf{events} in the plain past, kept in order with \textbf{\mr{आधी}}, \textbf{\mr{मग}}, \textbf{\mr{त्यानंतर}} and \textbf{\mr{शेवटी}}, and joined with \textbf{-\mr{ऊन}}. It has a \textbf{reason}, with \textbf{\mr{म्हणून}} or \textbf{\mr{कारण}}, and a \textbf{feeling}, with \textbf{\mr{मला} … \mr{आला}} or \textbf{\mr{वाटलं}}.

A good one ends by looking at what might have been: \textbf{\mr{जर} …, \mr{तर} \mr{मी} … \mr{गेलो} \mr{असतो}}.

\subsection*{Writing — connected composition}
On paper, in Devanagari, write \textbf{six lines} to your reader about a trip, real or invented, one line for each of these:

\begin{enumerate}
\raggedright
  \item the scene: which season it was, with \textbf{\mr{ऋतू}}, and what you were doing as it began, with \textbf{-\mr{त} \mr{होतो}} or \textbf{… \mr{तेव्हा}};
  \item the first thing that happened, with \textbf{\mr{आधी}}, and buying a ticket, with \textbf{\mr{तिकीट}};
  \item what happened next, with \textbf{\mr{मग}} or a \textbf{-\mr{ऊन}} form, and one thing that went wrong: \textbf{\mr{हरवली}} or \textbf{\mr{चुकली}};
  \item why, with \textbf{\mr{म्हणून}}, \textbf{\mr{कारण}} or \textbf{-\mr{मुळे}};
  \item how it ended and how you felt, with \textbf{\mr{शेवटी}} and \textbf{\mr{मला} … \mr{आला}}, \textbf{\mr{वाटलं}} or \textbf{\mr{झालं}};
  \item what might have been, with \textbf{… \mr{असतो}}, and one thing the trip taught you, with \textbf{\mr{ज्ञान}}.
\end{enumerate}

There is no model in this lesson. Use only words you already write. The work is choosing the right past for each line, and joining them.

\begin{tcolorbox}[breakable,colback=teal!4,colframe=teal!35!black,title={Before you move on}]
Now read your lines as you read Arun's monsoon day. Mark every verb \textbf{s} for scene, \textbf{e} for event. Line one should be mostly \textbf{s}; lines two and three mostly \textbf{e}.

Then three checks. A verb with an object in the plain past follows the thing: \textbf{-\mr{ने}} on he, she or a name. \textbf{\mr{ऋतू}} starts with the letter \textbf{\mr{ऋ}}, not \textbf{\mr{र}}. \textbf{\mr{ज्ञान}} starts with \textbf{\mr{ज्ञ}}, with \textbf{\mr{ञ}} inside it.

Last, read line six alone. Do both halves of the if wear \textbf{\mr{असता}} or \textbf{\mr{असतो}}?
\end{tcolorbox}
\section[(review)]{Review — it depends, and what would have been}

\label{lesson:MR-R287-depends-and-would}



\begin{quote}
Say \textquotedblleft{}leave early, otherwise you'll miss the train\textquotedblright{}. (\emph{lavkar nighā, nāhītar gāḍī cukel}.) Then \textquotedblleft{}I like the monsoon season most\textquotedblright{}. (\emph{malā pāvsāḷā hā ṛtū sarvāt jāsta āvaḍto}.) Then read aloud a line of yours with \textbf{\mr{ज्ञान}}.
\end{quote}

\subsection*{Your turn}
Say these aloud:

\begin{itemize}
\raggedright
  \item Will we go tomorrow? It depends on the weather.
  \item It depends on you.
  \item If I had had time, I would have come.
  \item If it had not rained, we would have gone to the river.
  \item one trip that depends on something, and one that would have gone better, if …
\end{itemize}

\begin{tcolorbox}[breakable,colback=teal!4,colframe=teal!35!black,title={Before you move on}]
\textquotedblleft{}It depends on\textquotedblright{}? (\textbf{… -\mr{वर} \mr{अवलंबून} \mr{आहे}}.) \textquotedblleft{}I would have gone\textquotedblright{}? (\textbf{\mr{मी} \mr{गेलो} \mr{असतो}}.) \textquotedblleft{}If I had had time\textquotedblright{}? (\textbf{\mr{जर} \mr{मला} \mr{वेळ} \mr{असता}}.) Which letters did \textbf{\mr{ऋतू}} and \textbf{\mr{ज्ञान}} bring? (\textbf{\mr{ऋ}} and \textbf{\mr{ञ}}.)
\end{tcolorbox}
